% Generated. Do not edit.
% Projection scriptural-date-location/evidence, version 1: the Scriptural Date and
% Location page, generated whole from research/chronology.toml, the Scripture
% chronology corpus and the source library by tools/tpt document-types refresh.
% document: liturgy/roman-rite/1962/propers/temporal/56-sixteenth-after-pentecost
\section*{Appendix: Evidence for the Scriptural Dates}
\label{sec:date-location-evidence}
{\small For each passage, as the Scriptural Date and Location table names it: what is dated, how the corpus weighs each answer, the source's own words, and the sources cited. A subject the corpus does not date is reported as such and given no date.\par}
{\small
\par\smallskip\noindent\textit{Offertory: Ps. 39:14--15 (Vg.; Heb. 40)}\par\nobreak
\begin{itemize}[leftmargin=1.2em,itemsep=1pt,topsep=1pt,parsep=0pt]
\item \textbf{Traditional attribution} (disputed). David's reign: B.C.~1055--1015 (``reigned from 1055 to 1015 B.C.''; The Catholic Encyclopedia, Volume 4 (New York, 1908), article 'David, King,' printed p. 642, opening paragraph). Place: Hebron, then Jerusalem (The Holy Bible, Douay-Rheims (Challoner revision), Project Gutenberg eBook 1581: 2 Kings 5:5).
\item \textbf{Historical setting}. No date: the chronology corpus has not yet dated this (research pending).
\item \textbf{Composition} (preferred). Psalms: Before c.~165 B.C. (``before the Maccabean period, around 165 B.C.; no individual psalm can be dated securely''; Catholic critical chronology; New American Bible, Revised Edition --- USCCB web state (2026), Psalms introduction).
\end{itemize}
\par\smallskip\noindent\textit{Communion: Ps. 70:16--18 (Vg.; Heb. 71)}\par\nobreak
\begin{itemize}[leftmargin=1.2em,itemsep=1pt,topsep=1pt,parsep=0pt]
\item \textbf{Traditional attribution}. As under Offertory: Ps. 39:14--15 (Vg.; Heb. 40).
\item \textbf{Superscription setting}. No date: the corpus holds a date for this that it reports from another edition by exception, which is not printed until the corpus records the source's own qualification of it (research pending).
\item \textbf{Historical setting} (preferred). David's flight from the face of Absalom: When he fled from the face of his son Absalom (The Holy Bible, Douay-Rheims (Challoner revision), Project Gutenberg eBook 1581: Ps. 3:1).
\item \textbf{Composition}. As under Offertory: Ps. 39:14--15 (Vg.; Heb. 40).
\end{itemize}
\par\smallskip\noindent\textit{Introit: Ps. 85:3, 5 (Vg.; Heb. 86), Ps. 85:1 (Vg.; Heb. 86)}\par\nobreak
\begin{itemize}[leftmargin=1.2em,itemsep=1pt,topsep=1pt,parsep=0pt]
\item \textbf{Traditional attribution}. As under Offertory: Ps. 39:14--15 (Vg.; Heb. 40).
\item \textbf{Historical setting}. As under Offertory: Ps. 39:14--15 (Vg.; Heb. 40).
\item \textbf{Composition}. As under Offertory: Ps. 39:14--15 (Vg.; Heb. 40).
\end{itemize}
\par\smallskip\noindent\textit{Alleluia: Ps. 97:1 (Vg.; Heb. 98)}\par\nobreak
\begin{itemize}[leftmargin=1.2em,itemsep=1pt,topsep=1pt,parsep=0pt]
\item \textbf{Traditional attribution}. As under Offertory: Ps. 39:14--15 (Vg.; Heb. 40).
\item \textbf{Historical setting}. As under Offertory: Ps. 39:14--15 (Vg.; Heb. 40).
\item \textbf{Composition}. As under Offertory: Ps. 39:14--15 (Vg.; Heb. 40).
\item \textbf{Prophetic referent} (disputed). The Nativity of Our Lord: B.C.~3 (``the year of Rome 750 which he styles 3 B.C.''; The Catholic Encyclopedia, Volume 3 (New York, 1908)); c.~B.C.~7 (``probably the year 7 B.C.''; The Catholic Encyclopedia, Volume 3 (New York, 1908)). Place: Bethlehem as the Christian referent (The Holy Bible, Douay-Rheims (Challoner revision), Project Gutenberg eBook 1581: Matt. 2:1).
\end{itemize}
\par\smallskip\noindent\textit{Gradual: Ps. 101:16--17 (Vg.; Heb. 102)}\par\nobreak
\begin{itemize}[leftmargin=1.2em,itemsep=1pt,topsep=1pt,parsep=0pt]
\item \textbf{Historical setting}. As under Offertory: Ps. 39:14--15 (Vg.; Heb. 40).
\item \textbf{Composition}. As under Offertory: Ps. 39:14--15 (Vg.; Heb. 40).
\end{itemize}
\par\smallskip\noindent\textit{Gospel: Luke 14:1--11}\par\nobreak
\begin{itemize}[leftmargin=1.2em,itemsep=1pt,topsep=1pt,parsep=0pt]
\item \textbf{Narrated event}. No date: the chronology corpus has not yet dated this (research pending).
\item \textbf{Composition} (disputed). The Gospel of St Luke: c.~A.D.~70 (``About the year 70''; The Catholic Encyclopedia, Volume 14 (New York, 1912)); Before the end of the Roman imprisonment, when the Acts was finished (De auctore, de tempore compositionis et de historica veritate Evangeliorum secundum Marcum et secundum Lucam (1912), responsa-i-ix).
\end{itemize}
\par\smallskip\noindent\textit{Epistle: Eph. 3:13--21}\par\nobreak
\begin{itemize}[leftmargin=1.2em,itemsep=1pt,topsep=1pt,parsep=0pt]
\item \textbf{Composition} (disputed). The Epistle to the Ephesians: A.D.~58--63 (``a period between 58 and 63''; The Catholic Encyclopedia, Volume 5 (New York, 1909)); A.D.~61 (``(Philemon; Colossians; Ephesians; Philippians), 61''; The Catholic Encyclopedia, Volume 11 (New York, 1911)). Place: Rome or Caesarea; recipients in Asia Minor (The Catholic Encyclopedia, Volume 5 (New York, 1909)).
\end{itemize}
\par}
